\documentclass[quick,pro]{debate}
\motion{语言的边界是 / 不是人类的边界}
\stance{正方}{是}
\begin{document}
\debatetitle
\begin{thesisbox}人比别的动物多走出去的每一步，都是语言铺的路。语言到不了的地方，一个人也许摸得到，人类走不到。\end{thesisbox}

\section{定义与判准（标准表述）}
\begin{fields}
\field{语言}{有词、有语法、用来表达和交流意思的符号系统。口语、文字、手语都是；音乐、绘画叫语言，只是比喻。}
\field{人类的边界}{人能想到、认识到、做到，并且能传给别人的范围的极限。}
\field{是}{两条边界基本重合：语言到不了的地方，人类在重要的方面也到不了。反例要稳定、成体系、对人重要。}
\field{判准}{语言到不了的地方，人类还能不能稳定地到达。\sub{我方要证明}{人超出本能的那一截以语言为条件；语言之外的，要么到不了，要么留不住。}\sub{对方要证明}{在语言之外，人能稳定地思考、创造，而且留得住。只说“有说不清的时候”不够。}}
\field{摸到与走到}{一个人摸到，是他的边界；说出来、传下去，人类才走到。}
\field{对方提偏向定义或判准时的 30 秒口径}{对方说语言只是嘴上的话：手语算不算语言？聋人有没有语言？对方说有一处不重合就不是：国界上也有飞过去的鸟，请回到“稳定、成体系、对人重要”。}
\end{fields}

\section{我方论点}
\begin{tabularx}{\linewidth}{@{}L{5.2em}Y Y L{6.5em}@{}}\toprule
\thead{标签}&\thead{一句话}&\thead{核心论据}&\thead{出处}\\\midrule
\textbf{没有词，走不到}&人超出动物的那一截思考，精确的数、别人心里想错了什么，跟着语言一起到&尼加拉瓜第一批 8 人（26.8 岁）错误信念 4 题平均对 0.5 题，第二批 10 人 3.8 题；“二十五年的社会经验也替代不了”。家庭手势者 4 人，1 到 3 全对，超过 3 精确答对不到一半&Pyers \& Senghas 2009；Spaepen 等 2011；维果茨基 1934\\
\textbf{说不出就接不上}&人类的边界是一代接一代往前走的边界；说不出的能守住，接不上力&石器实验 184 人：模仿没提高，比划翻一倍，说话翻四倍；奥杜威七十多万年几乎没变（作者推测）。轮扁：古之人与其不可传也死矣。性骚扰一词 1975 年定名&Morgan 等 2015；《庄子·天道》；History.com 2018\\
\bottomrule\end{tabularx}

\section{对方论点 → 一句话反驳}
\begin{tabularx}{\linewidth}{@{}Y Y Y@{}}\toprule
\thead{对方会说}&\thead{我方一句话回}&\thead{对方追问时}\\\midrule
失语的人照样算数、作曲&先认能。楼盖好了拆脚手架，楼还在；他们都是先有语言的&“聋孩子也能学数学”→ 同一篇写有些推理会慢，是后来接触了数字才学的\\
我们知道的比说的多（棋感、手艺）&一个人知道的比说的多，人类知道的不比说的多&“手艺传了几千年”→ 说不出的那点每代从零摸，这就是边界\\
爱因斯坦先有图像后找词&他先摸到；找到词、写成论文，人类才走到&“那边界在他脑子里”→ 在他脑子里的是他的边界\\
石器没语言也传了七十万年&传了，也停了七十多万年；守住不是往前走&“比划也翻一倍”→ 一倍守住，四倍往前\\
尼加拉瓜手语是孩子自己造的&造出“想”“知道”之前，他们走不到；人造词，词带人走&“那是人先到”→ 没有哪一步是人不带词先走过去的\\
5.6 说的是“我的”，维特根斯坦承认不可说&同意，所以我方论证人类，用的是实验；他让我们对不可说的沉默&“出处不支持你”→ 辩题问事实，不问出处\\
\bottomrule\end{tabularx}

\section{必问与必答}
\begin{points}{必问（对方不答就点名、重复、不换题）}
\item 说不出的东西，怎么让下一个人接着往前走？（全场问题）
\item 一辈子没有语言的人，您方有能精确数过三的例子吗？
\item 师傅走了，徒弟没摸到，人类还有没有这一点？
\end{points}
\begin{points}{必答（15 秒正面回答）}
\item 失语的人还能不能思考？→ 能。思考是有语言的几十年里搭起来的；他想的都是语言说得清的事。
\item 你方是不是说没有语言就不能思考？→ 不是。底子人人有，动物也有；我方说的是超出动物的那一截。
\item 你方是不是把个人开除出人类？→ 不是。个人摸到的，是人类的起点；说出来、传下去，才是人类走到。
\end{points}

\section{对辩战场概要（本赛制没有自由辩，三场对辩就是交锋主场）}
\begin{tabularx}{\linewidth}{@{}L{6em}Y Y Y@{}}\toprule
\thead{战场}&\thead{开场问题}&\thead{目标}&\thead{收束语}\\\midrule
三辩对辩：思考&一辈子没语言、能精确数到十的人，有吗？&失语者是先有语言的；婴儿的一加一是动物也有的底子&语言外面有底子，没有人类的边界\\
一辩对辩：摸到与走到&这个人走了以后，人类还剩下什么？&把全场分歧钉成“一个人摸到还是人类走到”&一个人摸到，是他的边界\\
二辩对辩：稳定&一个说不清的瞬间，算不算稳定？&守住中立判准的“稳定地到达、留得住”&守得住，接不上，就不算走到\\
\bottomrule\end{tabularx}

\section{如果……就……}
\begin{contingency}
\ifthen{对方定义语言为“嘴上的话”}{问手语算不算语言，回到中立定义。}
\ifthen{“没有词，走不到”被打穿}{收缩到“说不出就接不上”，用判准的“留得住”解释为什么仍然占优。}
\end{contingency}

分工：一辩守定义判准、为结辩取材，二辩守论点一，三辩守论点二。

\begin{recordcard}
\record{反一立论结束}{对方立论的主干}{\opt{对方立论建立在失语者还能思考上}\opt{对方立论建立在默会知识上}\opt{对方立论建立在语言只是交流工具上}\opt{对方立论的主干不在以上几条时}}{正二申论·拆反方立论}
\record{反二申论结束}{反二申论主打什么}{\opt{反二申论主打失语者还能思考}\opt{反二申论主打爱因斯坦先想后找词}\opt{反二申论主打石器没有语言也传了七十万年}\opt{反二申论不在以上几条时}}{正二申论·回应反二}
\record{一辩对辩结束}{对方全场最有力的一点}{\opt{对方全场最有力的是失语者还能思考}\opt{对方全场最有力的是默会知识}\opt{对方全场最有力的是石器没有语言也传了七十万年}\opt{对方最有力的一点不在以上几条时}}{正三申论·处理反方最强点}
\record{二辩对辩结束}{全场主战场落在}{\opt{全场主战场落在失语者能不能思考}\opt{全场主战场落在说不出的能不能传下去}\opt{主战场不清楚时}}{正一结辩·归纳分歧}
\record{二辩对辩结束}{反方全场最有力的一点}{\opt{反方全场最有力的是失语者还能思考}\opt{反方全场最有力的是默会知识}\opt{反方最有力的一点不在以上几条时}}{正一结辩·处理最强点}
\record{反一结辩结束}{反一收在哪里}{\opt{反一结辩收在语言只是工具不是边界}\opt{反一结辩收在说不出话的人还算不算完整的人}\opt{反一的收束不在以上几条时}}{正一结辩·回应反一（再接 40 字临场）}
\record{全场}{对方回避了哪些问题}{（写下来）}{申论与结辩里点名}
\end{recordcard}
\end{document}
